%!TEX root=../../../note
\subsection{Limits}
A sequence of real numbers is a function defined on the set $N = \set{1, 2, 3, \cdots}$ of natural numbers, whose range is contained 
in the set $\R$ of real numbers. 
\begin{equation*}
    X: \N \to \R
\end{equation*}
\begin{equation*}
    x_n: \text{ element of the sequence. }
\end{equation*}
Notation: $\set{x_n }, \set{x_n }_{n = 1}^{\infty}$

\begin{definition}
    A sequence $\set{x_n }$ in $\R$ is said to converge to $a \in \R$, or $a$ is said to be a limit of $\set{x_n }$ 
    , if for every $\epsilon > 0$ there exists a natural number $n_\epsilon \in \N$ such that for all $n \geq n_\epsilon$
    the term $x_n$ satisfy 
    \begin{equation*}
        \abs{x_n - a} < \epsilon. 
    \end{equation*} 

    The notation is 
    \begin{equation*}
        \lim_{n \to \infty} x_n = a. 
    \end{equation*}
\end{definition}

\begin{remark}[Reading the definition]
    In plain language: for every tolerance $\epsilon > 0$, all sufficiently late terms of the sequence lie inside the interval
    $\paren{a - \epsilon, a + \epsilon}$. Here $\epsilon$ is the tolerance, or precision level, we demand of the approximation
    $x_n \approx a$; the index $n_\epsilon$ tells us how late ``sufficiently late'' means, and in general depends on $\epsilon$.
    The phrase ``for all $n \geq n_\epsilon$'' says that once the sequence enters the $\epsilon$-band at index $n_\epsilon$, it
    stays inside that band for every later index -- it does not escape and re-enter. Since $n_\epsilon$ depends on $\epsilon$,
    tightening the tolerance (taking $\epsilon$ smaller) typically forces $n_\epsilon$ to be larger: the sequence may need to
    wait longer before it settles into a narrower band around $a$.
\end{remark}

The same ``eventually close'' idea appears in the $\epsilon$--$\delta$ definition of a function limit. Consider the oscillating function
\[
    f(x)=x\sin\paren{\frac1x}\qquad (x\ne0).
\]
Its oscillations become more frequent near $0$, but their amplitude is at most $|x|$, so the graph is squeezed toward the limiting value $0$.

\begin{center}
\begin{tikzpicture}[x=2.4cm,y=2.2cm,>=Stealth]
    % The epsilon-band and the delta-neighbourhood.
    \fill[red!8] (-0.9,-0.35) rectangle (0.9,0.35);
    \draw[densely dashed,red!75!black] (-0.9,0.35) -- (0.9,0.35) node[right] {$\epsilon$};
    \draw[densely dashed,red!75!black] (-0.9,-0.35) -- (0.9,-0.35) node[right] {$-\epsilon$};
    \draw[densely dashed,blue!70!black] (-0.35,-0.85) -- (-0.35,0.85) node[above] {$a-\delta$};
    \draw[densely dashed,blue!70!black] (0.35,-0.85) -- (0.35,0.85) node[above] {$a+\delta$};

    % Axes and the graph of x sin(1/x).
    \draw[->] (-0.95,0) -- (0.98,0) node[right] {$x$};
    \draw[->] (0,-0.9) -- (0,0.93) node[above] {$y$};
    \draw[very thick,teal!70!black,domain=0.025:0.9,samples=600,smooth]
        plot (\x,{\x*sin((1/\x) r)});
    \draw[very thick,teal!70!black,domain=-0.9:-0.025,samples=600,smooth]
        plot (\x,{\x*sin((1/\x) r)});
    \fill (0,0) circle (0.7pt);
    \node[below] at (0,-0.12) {$a=L=0$};
    \node[teal!70!black,anchor=south west] at (0.47,0.67) {$f(x)=x\sin(1/x)$};
    \node[blue!70!black,below] at (0,-0.88) {$|x-a|<\delta$};
\end{tikzpicture}
\end{center}

\begin{remark}[Proof idea]
    We want to control $|f(x)-L| = \abs{x\sin\paren{1/x}}$. Since $|\sin(\theta)| \leq 1$ for every real $\theta$, the factor
    $\sin(1/x)$ never contributes more than a factor of $1$, so $|f(x)| \leq |x|$: the size of $f(x)$ is bounded directly by
    the size of $x$ itself. This means we can make $|f(x)-L|$ small simply by making $|x-a|=|x|$ small, with no loss in the
    bound. Choosing $\delta=\epsilon$ then closes the chain of inequalities exactly, with no slack to spare.
\end{remark}

For a given $\epsilon>0$, choose $\delta=\epsilon$. If
\[
    0<|x-a|=|x|<\delta,
\]
then
\[
    |f(x)-L|=\abs{x\sin\paren{\frac1x}}
    \leq |x| \qquad \text{(since } |\sin(\theta)| \leq 1 \text{ for all } \theta \in \R \text{)}
    <\delta=\epsilon.
\]
Thus every $x$ in the $\delta$-neighbourhood of $a=0$ (apart from $a$ itself) has its function value inside the $\epsilon$-band around $L=0$.
